\documentclass[a4paper]{article}

%% Language and font encodings
\usepackage[english]{babel}
\usepackage[utf8]{inputenc}
\usepackage[T1]{fontenc}
\usepackage{amsfonts}

%% Sets page size and margins
\usepackage[a4paper,top=3cm,bottom=2cm,left=3cm,right=3cm,marginparwidth=1.75cm]{geometry}

%% Useful packages
\usepackage{amsmath}
\usepackage{amssymb}
\usepackage{amsthm}
\usepackage{graphicx}
\usepackage[section]{placeins}
\usepackage[colorinlistoftodos]{todonotes}
\usepackage[colorlinks=true, allcolors=blue]{hyperref}

\usepackage{titlesec}
\newcommand{\sectionbreak}{\clearpage}

%% Table packages
\usepackage{array}
\newcolumntype{P}[1]{>{\centering\arraybackslash}p{#1}}
\usepackage{multirow}

\title{Mathematical and Physical Modeling of Entangled Photon Propagation in Optical Fibers for MATLAB Simulation}
\author{Marco Cioci\\
Politecnico di Milano\\
Student ID: 10694772\\
\texttt{marco.cioci@mail.polimi.it}}
\date{\today}

\setlength{\parindent}{0pt}

\begin{document}
\maketitle

\tableofcontents

%========================================================
\newpage
\section{General Framework}

\subsection{Purpose and Scope}
The purpose of this document is to define a coherent working framework for the mathematical modeling and numerical simulation of polarization-entangled photon pairs propagating through optical fibers.

This text is not intended to appear directly in the final thesis manuscript. Instead, it serves as an internal technical guide throughout the development of the project. Its role is to support communication, maintain consistency of notation and assumptions, and provide a stable reference for the progressive validation of both the theoretical model and the MATLAB implementation.

The central objective is to organize the physical assumptions, mathematical formalism, and simulation strategy into a reproducible structure, so that each computational step can be traced back to a precise theoretical element. In this way, the simulator is not developed as an isolated numerical tool, but as a direct operational translation of the underlying quantum model.

\paragraph{Simulation notes.}
At the MATLAB level, this document should be treated as the reference specification for notation, model hierarchy, and implementation priorities. The simulator should therefore be organized into clearly identifiable modules, each corresponding to a specific theoretical block. This modular structure is essential both for step-by-step validation of intermediate results and for future extensions of the model.

\subsection{Theoretical Formalism and Notation}

The system is described within the framework of quantum information theory. The relevant mathematical structures include finite-dimensional Hilbert spaces, quantum states, tensor products, unitary operators, density operators, and projective measurements. Elements of classical optics are introduced only when required to justify the physical origin of the effective quantum transformations.

The notation and formal tools are aligned, as far as possible, with the framework of \textit{Quantum Computation and Quantum Information} by Nielsen and Chuang.

In this formalism:
\begin{itemize}
    \item single subsystems are represented by vectors in \(\mathbb{C}^2\),
    \item composite systems are described by tensor-product spaces,
    \item evolutions are represented by linear operators (unitary in the ideal case),
    \item measurements are described through Hermitian operators.
\end{itemize}

All states and operators are represented as vectors and matrices over \(\mathbb{C}\), with dimensions determined by the tensor-product structure of the system.

\paragraph{Simulation notes.}
The simulator adopts a quantum-information-oriented representation from the outset. Pure states are represented as complex vectors, mixed states as density matrices, and evolutions and measurements as matrix operators.

This abstract structure defines the numerical backbone of the simulator and is independent of the specific physical encoding used.


\subsection{Physical Scenario and Experimental Setting}

The physical setting considered in this work is the following. A spontaneous parametric down-conversion (SPDC) source generates pairs of polarization-entangled photons, ideally prepared in a Bell state. The two photons are injected into distinct optical paths, denoted \(A\) and \(B\), and propagate through independent optical fibers.

Each fiber acts on the polarization degree of freedom of the corresponding photon. In a general description, this action can be represented by a unitary operator acting on a single-qubit Hilbert space. However, motivated by experimental constraints and by the limited degree of control available on the full polarization transformation, the model is initially restricted to a simplified effective description.

After compensation of global polarization rotations, the dominant residual effect can be modeled as a relative phase shift between the horizontal and vertical polarization components. In the qubit representation, this corresponds to a rotation around the computational \(Z\)-axis of the Bloch sphere.

The objective at this stage is therefore not to reproduce the full electromagnetic propagation in optical fibers, but to construct an effective quantum model that captures the relevant impact of fiber propagation on entanglement and measurable correlations.

\paragraph{Simulation notes.}
At the simulator level, the physical scenario is reduced to the minimal set of elements required for computation: an initial bipartite quantum state, two labeled channels \(A\) and \(B\), and one local transformation acting on each subsystem.

In the first modeling stage, unnecessary optical detail is neglected and each fiber is described exclusively through its effective action on polarization. This leads to a local unitary representation of the form
\[
U_A \otimes U_B,
\]
where \(U_A, U_B \in \mathbb{C}^{2 \times 2}\) are unitary operators.

This abstraction defines the fundamental interface between the physical system and the simulator, and will later allow the introduction of more refined models, including stochastic phase evolution and non-unitary effects.
%========================================================
\section{Modeling Strategy}

\subsection{Progressive Modeling Approach}

The modeling approach follows a progressive refinement strategy, in which the system is described through successive levels of approximation of increasing physical complexity.

Each stage introduces additional physical effects while remaining consistent with the previous ones. This ensures analytical transparency and preserves continuity between simplified and refined models.

From a theoretical perspective, this approach isolates the contribution of individual mechanisms. From a computational perspective, it enables a layered simulator structure, where each refinement corresponds to a well-defined extension of the existing codebase.

%--------------------------------------------------------

\subsection{Hierarchy of Modeling Levels}

The modeling workflow is structured as a sequence of progressively refined descriptions:
\[
\begin{array}{c}
\text{Ideal Bell State} \\
\downarrow \\
\text{Local Phase Model (Z-rotations)} \\
\downarrow \\
\text{General Local Unitary Transformations} \\
\downarrow \\
\text{Mixed-State Description (Density Operators)} \\
\downarrow \\
\text{Noisy Quantum Channels} \\
\downarrow \\
\text{Effective Fiber Models}
\end{array}
\]

More explicitly:

\begin{itemize}
    \item \textbf{Level 1: Minimal phase model} \\
    Pure, maximally entangled state. Fiber action reduced to a relative phase. Analytical dependence of correlations on a single parameter.

    \item \textbf{Level 2: General unitary dynamics} \\
    Arbitrary polarization rotations described by local operators \(U_A \otimes U_B\).

    \item \textbf{Level 3: Mixed-state description} \\
    Density operators used to represent statistical mixtures and loss of coherence.

    \item \textbf{Level 4: Noisy quantum channels} \\
    Effective completely positive trace-preserving (CPTP) maps modeling depolarization and decoherence.

    \item \textbf{Level 5: Effective fiber models} \\
    Physically motivated descriptions, such as concatenations of random unitary segments or stochastic polarization evolution.
\end{itemize}

This hierarchy provides both a conceptual and an implementation roadmap. The initial levels remain analytically tractable and allow direct validation, while the later levels increase physical realism and connect to experimentally relevant effects.

\paragraph{Simulation notes.}
This hierarchy should be mirrored explicitly in the codebase. Each level should correspond to a distinct layer of functionality: state preparation, local unitary evolution, density-matrix handling, channel application, and advanced propagation models. This structure ensures that each refinement appears as an additional computational step rather than a structural modification of the simulator.

%--------------------------------------------------------

\subsection{Baseline Model: Pure-State Description}

At the first stage, the bipartite system is described as a pure state in \(\mathbb{C}^2 \otimes \mathbb{C}^2\), and fiber propagation is modeled through local unitary transformations. Particular emphasis is placed on relative phase shifts induced by the fibers.

The main objective is to analyze how these phase shifts affect measurable two-photon correlations. The relevant observables are operators of the form
\[
\sigma_i \otimes \sigma_j,
\qquad i,j \in \{x,y,z\},
\]
as well as operators corresponding to arbitrary measurement directions on the Bloch sphere.

This baseline model constitutes the simplest regime in which entanglement is preserved and can be probed exclusively through correlation observables. It provides a controlled setting for deriving analytical expressions and validating the numerical implementation.

\paragraph{Simulation notes.}
At this level, the simulator operates on pure two-qubit states represented as \(4 \times 1\) complex vectors. The core tasks are: generation of the Bell state, application of local operators \(U_A\) and \(U_B\), and evaluation of expectation values.

In this regime, the simulator enables direct comparison between analytical predictions and numerical results for correlation functions, establishing a first level of validation that will be extended in Chapter~4 through quantitative state characterization.

%--------------------------------------------------------
\subsection{Extension to Mixed-State Formalism}

At the next level, the model incorporates statistical effects and decoherence. In this regime, the pure-state description is no longer sufficient and must be replaced by a density-operator formalism.

This extension enables the description of:
\begin{itemize}
    \item reduced states of individual photons,
    \item loss of coherence due to unresolved degrees of freedom,
    \item effective non-unitary dynamics arising from environmental interactions or statistical averaging.
\end{itemize}

Within this framework, physically relevant quantities such as purity, fidelity, and concurrence can be defined and computed, allowing a quantitative characterization of mixedness and entanglement (see Chapter~4).

An important derived quantity is the reduced density operator of a subsystem, obtained via partial trace. For instance, the reduced state of subsystem \(A\) is
\[
\rho_A = \mathrm{Tr}_B(\rho_{AB}).
\]

For maximally entangled states such as \( |\Psi^+\rangle \), the reduced density operators are maximally mixed:
\[
\rho_A = \rho_B = \frac{1}{2} I.
\]

This property reflects the absence of local information despite strong nonlocal correlations, and provides the conceptual link between entanglement and mixedness that will be explored quantitatively in later sections.

Expectation values in the density-operator formalism are computed as
\[
\langle O \rangle = \mathrm{Tr}(\rho O),
\]
which generalizes the pure-state expression \( \langle \psi | O | \psi \rangle \).

\paragraph{Simulation notes.}
The simulator must be extended so that all routines operate on density matrices, including \(4 \times 4\) representations, partial trace operations, and observable evaluation.

This generalized framework enables the transition from deterministic evolution to ensemble-based descriptions, and supports the computation of state metrics used for validation and physical interpretation.

The pure-state implementation is preserved as a special case, ensuring consistency across all modeling levels.
%--------------------------------------------------------

\subsection{Simulator Design and Architecture}

The simulator is designed as a modular implementation of the theoretical model.

The initial modules include:
\begin{itemize}
    \item Bell-state generation,
    \item local unitary propagation,
    \item correlation evaluation.
\end{itemize}

More advanced modules introduce:
\begin{itemize}
    \item quantum channels,
    \item reduced density operators,
    \item entanglement measures,
    \item experimentally relevant observables such as Bell parameters and Stokes-like quantities.
\end{itemize}

The theoretical structure is directly reflected in the computational architecture, ensuring consistency, traceability, and extensibility.

The implementation distinguishes clearly between reusable computational kernels and higher-level experiment scripts, allowing the simulator to evolve into a structured and reusable research tool.

\paragraph{Simulation notes.}
At the current stage, this modular design is reflected in the directory structure:
\begin{itemize}
    \item \texttt{Core}: state vectors, operators, expectation values, correlation routines;
    \item \texttt{Models}: physical propagation models;
    \item \texttt{Utils}: validation routines and auxiliary tools;
    \item \texttt{Experiments}: scripts implementing complete simulation studies.
\end{itemize}

This separation enforces a strict distinction between general mathematical operations and physically specific modeling choices.

%========================================================
\section{Physical and Mathematical Representation}

\subsection{Quantum System and Hilbert Space Structure}

The physical system considered consists of polarization-entangled photon pairs generated by spontaneous parametric down-conversion (SPDC). The two photons propagate along distinct paths, denoted \(A\) and \(B\), forming a bipartite system.

In this work, the qubit is encoded in the polarization degree of freedom:
\[
|H\rangle \equiv |0\rangle,
\qquad
|V\rangle \equiv |1\rangle.
\]

Accordingly, the computational basis of the bipartite system is
\[
\{
|00\rangle,\,
|01\rangle,\,
|10\rangle,\,
|11\rangle
\},
\]
corresponding to
\[
\{
|H\rangle_A |H\rangle_B,\,
|H\rangle_A |V\rangle_B,\,
|V\rangle_A |H\rangle_B,\,
|V\rangle_A |V\rangle_B
\}.
\]

The Hilbert space of the total system is therefore
\[
\mathcal{H}_{AB} = \mathcal{H}_A \otimes \mathcal{H}_B
\cong \mathbb{C}^2 \otimes \mathbb{C}^2,
\]
which has dimension \(4\).

As a consequence:
\begin{itemize}
    \item pure states are represented by complex vectors of length \(4\),
    \item operators acting on the full system are \(4 \times 4\) matrices.
\end{itemize}

This construction specifies the concrete realization of the abstract quantum-information framework introduced previously.

\paragraph{Simulation notes.}
At MATLAB level, bipartite pure states are represented as \(4 \times 1\) complex vectors and operators as \(4 \times 4\) matrices. The basis ordering is fixed as
\[
\{|00\rangle,|01\rangle,|10\rangle,|11\rangle\},
\]
and must remain invariant, since all tensor products, observables, and density matrices depend on this convention.

This ordering is already enforced in the current codebase and must be preserved in all future extensions.

%--------------------------------------------------------

\subsection{Bell-State Initialization}

As a reference initial condition, we consider the Bell state
\[
|\Psi^{+}\rangle
=
\frac{1}{\sqrt{2}}
\left(
|H\rangle_A |V\rangle_B
+
|V\rangle_A |H\rangle_B
\right),
\]
which, in computational notation, is
\[
|\Psi^{+}\rangle
=
\frac{1}{\sqrt{2}}
\left(
|01\rangle + |10\rangle
\right).
\]

In the ordered computational basis
\[
\{|00\rangle, |01\rangle, |10\rangle, |11\rangle\},
\]
this state is represented by the column vector
\[
|\Psi^{+}\rangle
=
\frac{1}{\sqrt{2}}
\begin{pmatrix}
0\\
1\\
1\\
0
\end{pmatrix}.
\]

This state is maximally entangled and provides the natural input for the baseline version of the simulator. It also serves as the reference state with respect to which later quantities such as fidelity and Bell-inequality violation may be evaluated.

The choice of a Bell state is especially convenient because it provides both a clean theoretical benchmark and a physically relevant entangled resource. It also allows the effect of fiber propagation to be interpreted directly in terms of coherence and two-photon correlation changes.

\paragraph{Simulation notes.}
The simulator should include a dedicated state-preparation routine returning the Bell-state vector in canonical basis ordering. This reference state is useful for normalization checks, debugging, comparison with analytical formulas, and later evaluation of propagated-state fidelity.

At the current stage, the simulator already includes a Bell-state generator returning the four standard Bell states:
\[
|\Psi^{\pm}\rangle
=
\frac{1}{\sqrt{2}}\left(|01\rangle \pm |10\rangle\right),
\qquad
|\Phi^{\pm}\rangle
=
\frac{1}{\sqrt{2}}\left(|00\rangle \pm |11\rangle\right),
\]
although the present baseline model uses \( |\Psi^+\rangle \) as the default reference input.

%--------------------------------------------------------

\subsection{Fiber Propagation as Local Unitary Evolution}

Each photon propagates through its own optical fiber. At the level of the polarization qubit, the action of each fiber is described by a linear operator acting locally on the corresponding single-photon Hilbert space. In the most general ideal lossless description, this operator is unitary. Therefore, if \(U_A\) and \(U_B\) denote the single-qubit operators associated with the two fibers, the evolution of the initial two-photon state is
\[
|\psi_{\mathrm{out}}\rangle
=
(U_A \otimes U_B)\, |\psi_{\mathrm{in}}\rangle.
\]

This is the natural bipartite extension of single-qubit evolution and follows directly from the tensor-product structure of composite quantum systems. It is the appropriate starting point for later refinements, including random unitary perturbations, concatenated fiber sections, and effective channel descriptions.

In this framework, the two fibers act independently on the two polarization subsystems, and the total evolution is obtained by tensor composition of the corresponding local transformations. This description is valid as long as the fibers do not induce any direct interaction between the two photons.

\paragraph{Simulation notes.}
The fundamental evolution primitive in MATLAB should be the tensor-product action
\[
\psi_{\mathrm{out}} = (U_A \otimes U_B)\psi_{\mathrm{in}}.
\]
It is convenient to isolate this in a dedicated routine, so that any later modification of the single-fiber model automatically propagates to the full bipartite simulation.

This abstraction is already present in the implementation, where local single-qubit matrices are combined into the global operator \(U_A \otimes U_B\) before being applied to the input state.

%--------------------------------------------------------

%========================================================
\subsection{Reduced Phase Model}

Although a general fiber transformation may involve arbitrary polarization rotations, birefringence, and mode-coupling effects, the first approximation adopted here is simpler and directly motivated by the experimental control scheme. After partial compensation of polarization drifts, the dominant residual effect is modeled as a relative phase shift between the horizontal and vertical components, corresponding to an effective rotation around the \(Z\)-axis of the Bloch sphere.

At the single-qubit level, this is represented by the diagonal unitary
\[
U(\theta)
=
\begin{pmatrix}
1 & 0\\
0 & e^{i\theta}
\end{pmatrix},
\]
up to an irrelevant global phase. If the two fibers induce phases \(\theta_A\) and \(\theta_B\), then the total transformation is
\[
U_A \otimes U_B
=
U(\theta_A)\otimes U(\theta_B).
\]

Applying this operator to the reference Bell state \( |\Psi^+\rangle \) gives
\[
|\psi(\theta_A,\theta_B)\rangle
=
\frac{1}{\sqrt{2}}
\left(
e^{i\theta_B}|01\rangle
+
e^{i\theta_A}|10\rangle
\right).
\]
By factoring out a global phase, which has no effect on measurement statistics, this state can be rewritten as
\[
|\psi(\theta)\rangle
=
\frac{1}{\sqrt{2}}
\left(
|01\rangle
+
e^{i\theta}|10\rangle
\right),
\qquad
\theta = \theta_A - \theta_B.
\]

\paragraph{Interpretation of the model.}
In this reduced description, the physically relevant fiber action is encoded in a single parameter: the relative phase \(\theta\) between the two components of the entangled state. This parameter depends on propagation conditions such as birefringence, path length, and environmental fluctuations.

This model isolates the dominant effect affecting two-photon correlations while preserving a pure-state description. As a consequence, it provides an analytically transparent framework in which entanglement is maintained and all modifications arise solely from phase-induced coherence changes.

\paragraph{Role in the modeling hierarchy.}
The reduced phase model plays a central role as the first nontrivial propagation model. It provides the natural starting point for the simulator, allowing one to derive explicit analytical predictions and to validate the numerical implementation in a controlled setting.

At the same time, it constitutes the reference point for subsequent extensions. More general unitary transformations, ensemble descriptions, and effective non-unitary evolutions can all be understood as refinements of this baseline model.

\paragraph{Operational use in the simulator.}
The first stage of the simulator is therefore built around the family of states
\[
|\psi(\theta)\rangle
=
\frac{1}{\sqrt{2}}
\left(
|01\rangle + e^{i\theta}|10\rangle
\right),
\]
from which one computes expectation values of local and nonlocal observables, correlation functions, and Bell-type quantities as functions of \(\theta\).

\paragraph{Simulation notes.}
At implementation level, the simulator should allow \(\theta_A\) and \(\theta_B\) as input parameters, while recognizing that all observable quantities depend only on the relative phase \(\theta=\theta_A-\theta_B\). In practice, both parametrizations are useful: the two-phase form for physical interpretation, and the reduced one-parameter form for analytical checks and parameter sweeps.

The first validated MATLAB study consists of a parameter sweep in \(\theta\), with numerical evaluation of state vectors, expectation values, and correlation quantities. This provides a minimal but complete simulation pipeline, forming the basis for all subsequent developments.
%--------------------------------------------------------

\subsection{Correlation Observables: Theory and Implementation}

At the present stage of development, the observables explicitly implemented in the simulator are the standard aligned two-qubit Pauli correlations:
\[
\sigma_x \otimes \sigma_x,\qquad
\sigma_y \otimes \sigma_y,\qquad
\sigma_z \otimes \sigma_z.
\]

These are constructed from the single-qubit Pauli matrices
\[
\sigma_x=
\begin{pmatrix}
0 & 1\\
1 & 0
\end{pmatrix},
\qquad
\sigma_y=
\begin{pmatrix}
0 & -i\\
i & 0
\end{pmatrix},
\qquad
\sigma_z=
\begin{pmatrix}
1 & 0\\
0 & -1
\end{pmatrix}.
\]

Given a pure two-qubit state \( |\psi\rangle \), the expectation values are defined as
\[
\langle \sigma_i \otimes \sigma_i \rangle
=
\langle \psi|
(\sigma_i \otimes \sigma_i)
|\psi\rangle,
\qquad i\in\{x,y,z\}.
\]

\paragraph{Analytical behavior under the reduced phase model.}

Within the reduced phase model, the initial Bell state
\[
|\Psi^+\rangle = \frac{|01\rangle + |10\rangle}{\sqrt{2}}
\]
evolves into
\[
|\psi(\theta)\rangle
=
\frac{|01\rangle + e^{i\theta}|10\rangle}{\sqrt{2}}.
\]

For this family of states, the correlation observables can be computed analytically:
\[
\langle \sigma_x \otimes \sigma_x \rangle = \cos\theta,
\qquad
\langle \sigma_y \otimes \sigma_y \rangle = \cos\theta,
\qquad
\langle \sigma_z \otimes \sigma_z \rangle = -1.
\]

These results show that the fiber-induced phase \(\theta\) affects only the transverse correlations in the \(x\)- and \(y\)-directions, while the longitudinal correlation along \(z\) remains invariant. This reflects the fact that the relative phase modifies coherence between computational-basis components without changing their populations.

\paragraph{Simulation notes.}
At implementation level, the simulator contains routines for:
\begin{itemize}
    \item generating the Pauli matrices,
    \item constructing tensor-product operators,
    \item evaluating expectation values,
    \item packaging correlation outputs into structured data containers.
\end{itemize}

Additional helper routines verify normalization of input state vectors and unitarity of operators used in the fiber model.

These observables constitute the primary diagnostic outputs of the simulator and provide the most direct bridge between analytical predictions and numerical results. Their dependence on \(\theta\) defines the first validation pipeline of the code.

%------------------------------------------------------------
%========================================================
\subsection{From Deterministic Evolution to Ensemble Description}

The reduced phase model introduced in the previous sections describes a deterministic unitary evolution of a pure entangled state, parametrized by a single phase \(\theta\).

In realistic fiber propagation, however, this phase is not constant but fluctuates due to environmental effects such as birefringence variations and thermal instability. As a consequence, each photon pair experiences a different phase shift.

When these fluctuations occur on timescales shorter than the measurement resolution, individual realizations cannot be resolved experimentally. The observed system is therefore not described by a single pure state, but by an ensemble of states corresponding to different phase realizations.

\medskip

In order to describe such situations, it is necessary to introduce the density operator formalism. A quantum state is represented not only by a state vector, but more generally by an operator \(\rho\), defined as
\[
\rho = \sum_k p_k \, |\psi_k\rangle\langle \psi_k|,
\]
where \(\{p_k\}\) is a probability distribution and \(|\psi_k\rangle\) are normalized pure states. This representation captures both classical statistical uncertainty and loss of coherence arising from unresolved degrees of freedom.

\medskip

In the present context, the effective description is obtained by averaging over the phase distribution:
\[
\rho = \mathbb{E}_{\theta}\left[ |\psi(\theta)\rangle\langle \psi(\theta)| \right].
\]

This transition does not introduce new dynamics at the level of individual realizations, which remain unitary, but reflects a loss of accessible phase information. The resulting state is said to be \emph{mixed}, even though each underlying realization is pure.

\medskip

The density operator \(\rho\) satisfies the following fundamental properties:
\[
\rho^\dagger = \rho, 
\qquad
\rho \ge 0, 
\qquad
\mathrm{Tr}(\rho) = 1.
\]
These conditions ensure that \(\rho\) provides a complete and physically consistent description of the quantum state.

\medskip

Expectation values of observables are computed in this formalism as
\[
\langle O \rangle = \mathrm{Tr}(\rho O),
\]
which generalizes the pure-state expression \(\langle \psi | O | \psi \rangle\).

\medskip

This framework forms the basis for the modeling of phase-noise effects and, more generally, for the description of non-deterministic or partially coherent quantum evolutions. In the following sections, it will be extended to include more general transformations acting directly on density operators, and in Chapter~4 it will be used to introduce quantitative measures of mixedness and entanglement.

%========================================================
\subsection{Quantum Channels and Effective Depolarization}

The ensemble description introduced in the previous section accounts for statistical fluctuations of the relative phase parameter by averaging over a family of pure states. In that setting, each individual realization is still described by a unitary transformation, and the resulting mixed state arises from classical uncertainty over these realizations.

However, in realistic fiber propagation, the loss of coherence may originate not only from classical fluctuations of a control parameter, but also from the coupling between polarization and additional physical degrees of freedom that are not directly observed, such as frequency or temporal modes. In such cases, the effective evolution of the polarization state can no longer be described as an average over unitary transformations acting on a closed system. Instead, it must be formulated as a more general transformation acting directly on the density operator.

\medskip

This motivates the introduction of the concept of a \emph{quantum channel}, which provides the appropriate mathematical framework for describing effective non-unitary evolution.

\paragraph{Quantum channels as effective evolution maps}

A quantum channel is defined as a linear map
\[
\mathcal{E} : \rho \mapsto \mathcal{E}(\rho),
\]
acting on density operators, such that it is completely positive and trace preserving. These properties ensure that the output state remains a valid physical density operator.

A general representation of such maps is given by the operator-sum (Kraus) form:
\[
\mathcal{E}(\rho)
=
\sum_{\mu} K_{\mu} \rho K_{\mu}^\dagger,
\qquad
\sum_{\mu} K_{\mu}^\dagger K_{\mu} = I,
\]
where the operators \(K_{\mu}\) encode the effective action of the physical process on the system.

\medskip

In the present context, the channel formalism provides a natural extension of the unitary model of fiber propagation. While ideal fibers are described by local unitary operators of the form \(U_A \otimes U_B\), realistic propagation may involve mechanisms that effectively couple polarization to unobserved degrees of freedom. The resulting polarization state is then obtained by tracing out these auxiliary variables, leading to a non-unitary evolution at the level of the reduced system.

\paragraph{Physical origin of effective depolarization}

One of the main mechanisms responsible for such non-unitary behavior in optical fibers is the frequency dependence of the polarization transformation. When the input photons are described by finite-duration wave packets, their spectral components may experience different polarization-dependent phase shifts.

As a consequence, distinct polarization components can acquire different group delays and become partially or completely separated in time. When the temporal separation exceeds the coherence time of the detection process, the interference between these components is lost. The system is therefore no longer described by a coherent superposition, but by a statistical mixture.

\medskip

From the point of view of the polarization degree of freedom alone, this loss of temporal overlap manifests as a loss of coherence, leading to an effective depolarization of the state. This provides a direct physical interpretation of the transition from a pure-state description to a mixed-state one.

\paragraph{Depolarizing channel model}

As a first effective model of this behavior, we introduce the depolarizing channel. In the bipartite two-qubit setting, its action on a density operator \(\rho\) can be written as
\[
\mathcal{E}(\rho)
=
(1 - p)\,\rho
+
\frac{p}{4}\, I_4,
\]
where \(I_4\) denotes the identity operator on the four-dimensional Hilbert space and \(p \in [0,1]\) is a parameter quantifying the strength of the depolarization.

\medskip

This model interpolates between two limiting cases:

\begin{itemize}
    \item for \(p = 0\), the evolution is trivial and the input state is preserved;
    \item for \(p = 1\), the output state is completely mixed, corresponding to a total loss of polarization information.
\end{itemize}

The depolarizing channel is isotropic, in the sense that it does not privilege any particular basis. For this reason, it provides a natural first approximation to situations in which polarization information is uniformly degraded.

\paragraph{Application to Bell states}

When the input state is the Bell state
\[
|\Psi^+\rangle = \frac{|01\rangle + |10\rangle}{\sqrt{2}},
\]
with associated density operator \(\rho_0 = |\Psi^+\rangle\langle\Psi^+|\), the action of the depolarizing channel produces the family of states
\[
\rho(p)
=
(1 - p)\,\rho_0
+
\frac{p}{4}\, I_4.
\]

These states provide a simple yet nontrivial model for the degradation of entanglement under fiber propagation. In particular, they allow one to study the transition from a maximally entangled pure state to a mixed state with reduced or vanishing entanglement.

\paragraph{Connection with state characterization}

The introduction of quantum channels completes the modeling of the dynamical evolution of the system. The resulting states are, in general, mixed and must therefore be analyzed using the density-operator formalism.

In the following chapter, quantitative measures such as purity, fidelity, and concurrence will be used to characterize the effects of depolarization on mixed bipartite states, providing a detailed assessment of coherence loss and entanglement degradation.



%========================================================
%========================================================
%========================================================
\section{Quantitative Characterization of Bipartite Quantum States}

\subsection{Reduced Density Operators}

In bipartite quantum systems, the state of an individual subsystem is generally described by a density operator obtained through a partial trace over the complementary subsystem. This formalism accounts for correlations that may be present at the global level.

Let the composite system be defined on the Hilbert space
\[
\mathcal{H}_{AB} = \mathcal{H}_A \otimes \mathcal{H}_B,
\]
with $\mathcal{H}_A \cong \mathbb{C}^2$ and $\mathcal{H}_B \cong \mathbb{C}^2$. A generic global state is described by a density operator
\[
\rho_{AB} \in \mathbb{C}^{4 \times 4}.
\]

The reduced density operators are defined as
\[
\rho_A = \mathrm{Tr}_B(\rho_{AB}), \qquad
\rho_B = \mathrm{Tr}_A(\rho_{AB}).
\]

These operators reproduce all expectation values of observables acting locally on the subsystem. For any observable $O_A$ acting on subsystem $A$,
\[
\langle O_A \rangle = \mathrm{Tr}(\rho_{AB}(O_A \otimes I_B)) = \mathrm{Tr}(\rho_A O_A).
\]

An explicit construction can be written as
\[
\rho_A = \sum_{j=0}^{1} (I_A \otimes \langle j|)\,\rho_{AB}\,(I_A \otimes |j\rangle),
\]
which shows that the partial trace eliminates the degrees of freedom associated with the unobserved subsystem.

For maximally entangled Bell states, such as
\[
|\Psi^+\rangle = \frac{|01\rangle + |10\rangle}{\sqrt{2}},
\]
one obtains
\[
\rho_A = \rho_B = \frac{I}{2},
\]
indicating that each subsystem is locally maximally mixed despite the global state being pure.

The same result holds for the phase-evolved state
\[
|\psi(\theta)\rangle = \frac{|01\rangle + e^{i\theta}|10\rangle}{\sqrt{2}},
\]
where the relative phase modifies global coherence but leaves local states unchanged:
\[
\rho_A(\theta) = \rho_B(\theta) = \frac{I}{2}.
\]

A detailed derivation of this result, including the explicit matrix representation and term-by-term partial trace, is provided in Appendix~A.

By linearity of the partial trace, this property is preserved under convex combinations. Thus, for the ensemble state
\[
\rho_{AB}^{(\mathrm{ens})}
=
\frac{1}{N}
\sum_{k=1}^{N}
|\psi(\theta_k)\rangle\langle\psi(\theta_k)|,
\]
one finds
\[
\rho_A^{(\mathrm{ens})} = \rho_B^{(\mathrm{ens})} = \frac{I}{2}.
\]

This leads to an important conceptual point: the same reduced density operator
\[
\rho_A = \frac{I}{2}
\]
can arise from physically distinct global states.

On one hand, it may describe a classical ensemble
\[
\rho^{(\mathrm{class})}
=
\frac{1}{2}|0\rangle\langle 0|
+
\frac{1}{2}|1\rangle\langle 1|,
\]
where the system is in a definite state but the observer lacks knowledge.

On the other hand, it may result from entanglement:
\[
\rho_A^{(\mathrm{ent})}
=
\mathrm{Tr}_B\big(|\Psi^+\rangle\langle\Psi^+|\big)
=
\frac{I}{2}.
\]

Since both cases yield identical expectation values for all local observables,
\[
\mathrm{Tr}(\rho^{(\mathrm{class})} O_A)
=
\mathrm{Tr}(\rho_A^{(\mathrm{ent})} O_A),
\]
they are indistinguishable by local measurements alone.

The distinction emerges only at the global level, where coherence and correlations must be analyzed through quantities such as purity and bipartite correlation functions (see Appendix~A for an explicit comparison).

Finally, the reduced density operator directly determines measurable local quantities. For Pauli observables,
\[
\langle \sigma_x \rangle = \mathrm{Tr}(\rho_A \sigma_x), \quad
\langle \sigma_y \rangle = \mathrm{Tr}(\rho_A \sigma_y), \quad
\langle \sigma_z \rangle = \mathrm{Tr}(\rho_A \sigma_z),
\]
and for $\rho_A = I/2$ all expectation values vanish, indicating complete local depolarization.

%-------------------------------------------------------------------
\subsection{Purity and Mixedness}

Having introduced reduced density operators in Section 4.1, we now require a quantitative criterion to distinguish whether a quantum state is pure or mixed. This distinction is essential for interpreting the physical effects of fiber propagation, especially when extending the model beyond deterministic unitary evolution.

A quantum state described by a density operator $\rho$ is said to be \emph{pure} if it can be written as
\[
\rho = |\psi\rangle \langle \psi|,
\]
for some normalized state vector $|\psi\rangle$. Otherwise, the state is said to be \emph{mixed}, meaning that it cannot be represented by a single state vector but only as a statistical ensemble
\[
\rho = \sum_k p_k |\psi_k\rangle \langle \psi_k|, \qquad p_k \ge 0, \quad \sum_k p_k = 1.
\]

A convenient and operational way to distinguish between these two cases is provided by the \emph{purity}, defined as
\[
\gamma(\rho) = \operatorname{Tr}(\rho^2).
\]

This quantity satisfies the fundamental properties
\[
\gamma(\rho) = 1 \quad \Longleftrightarrow \quad \rho \text{ is pure},
\]
\[
\gamma(\rho) < 1 \quad \Longleftrightarrow \quad \rho \text{ is mixed}.
\]

This follows directly from the spectral decomposition of $\rho$. If
\[
\rho = \sum_i \lambda_i |i\rangle \langle i|,
\]
with $\lambda_i \ge 0$ and $\sum_i \lambda_i = 1$, then
\[
\gamma(\rho) = \sum_i \lambda_i^2.
\]
The purity reaches its maximum value $1$ only when a single eigenvalue is equal to $1$ and all others vanish, corresponding to a pure state. Otherwise, the sum of squares is strictly smaller than $1$.

More generally, for a system with Hilbert space dimension $d$, the purity is bounded by
\[
\frac{1}{d} \le \gamma(\rho) \le 1,
\]
where the lower bound is attained by the maximally mixed state $\rho = I/d$.

\medskip

For single-qubit systems, purity admits a particularly transparent interpretation in terms of the Bloch vector representation introduced in Section 4.1. Any qubit density operator can be written as
\[
\rho = \frac{1}{2}\left(I + \vec{r} \cdot \vec{\sigma}\right),
\]
where $\vec{r} \in \mathbb{R}^3$ is the Bloch vector and $\vec{\sigma} = (\sigma_x, \sigma_y, \sigma_z)$ is the vector of Pauli matrices. In this representation, one finds
\[
\gamma(\rho) = \frac{1 + \|\vec{r}\|^2}{2}.
\]
Therefore,
\[
\|\vec{r}\| = 1 \;\Rightarrow\; \rho \text{ is pure},
\qquad
\|\vec{r}\| = 0 \;\Rightarrow\; \rho = \frac{I}{2} \text{ is maximally mixed}.
\]

This result provides a direct connection between the Bloch/Stokes parameters computed in the simulator and the degree of mixedness of the state.

\medskip

As already emphasized in Section 4.1, mixedness does not necessarily imply classical noise. A state
may be mixed either because it belongs to a classical statistical ensemble or because it arises as
the reduced state of a larger entangled system.

In the present context, both mechanisms are relevant but play distinct roles. Mixedness induced by
entanglement is intrinsic to the bipartite structure of the system and is already present for pure
global states. By contrast, ensemble-induced mixedness arises from averaging over random parameters,
such as phase fluctuations in fiber propagation, and leads to a genuine loss of global coherence.

To illustrate the first mechanism, we consider a paradigmatic example provided by the Bell state
\[
|\Psi^+\rangle = \frac{|01\rangle + |10\rangle}{\sqrt{2}}.
\]
Its global density operator
\[
\rho_{AB} = |\Psi^+\rangle \langle \Psi^+|
\]
is pure, and therefore satisfies
\[
\gamma(\rho_{AB}) = 1.
\]
However, its reduced density operators are
\[
\rho_A = \rho_B = \frac{I}{2},
\]
which are maximally mixed and satisfy
\[
\gamma(\rho_A) = \gamma(\rho_B) = \frac{1}{2}.
\]

Thus, a globally pure entangled state may generate locally mixed subsystem states. This distinction between global purity and local mixedness is central to the interpretation of bipartite entangled systems.

\medskip

In the context of the present model, this distinction plays a crucial role. For the phase-evolved Bell-like state
\[
|\psi(\theta)\rangle = \frac{|01\rangle + e^{i\theta}|10\rangle}{\sqrt{2}},
\]
the evolution is generated by local unitary transformations, and therefore the global state remains pure for all values of $\theta$:
\[
\gamma(\rho(\theta)) = 1.
\]
At the same time, as shown in Section 4.1, the reduced states remain
\[
\rho_A(\theta) = \rho_B(\theta) = \frac{I}{2},
\]
so that their purity is constant and minimal:
\[
\gamma(\rho_A(\theta)) = \gamma(\rho_B(\theta)) = \frac{1}{2}.
\]

Therefore, the introduction of a relative phase modifies the correlation observables without affecting either global purity or local mixedness.

\medskip

The notion of purity becomes particularly relevant when extending the model to include more realistic physical effects. 
While deterministic unitary evolution preserves purity, ensemble averaging over random parameters or the inclusion of frequency-dependent propagation effects, such as polarization mode dispersion, can lead to genuine mixed states. 

In particular, for the random-phase ensemble introduced above, the global density operator depends on the ensemble coherence
\[
\mu = \langle e^{i\theta} \rangle,
\]
and its purity is given by
\[
\gamma\!\left(\rho_{AB}^{(\mathrm{ens})}\right)
=
\frac{1 + |\mu|^2}{2}.
\]
Equivalently,
\[
\gamma\!\left(\rho_{AB}^{(\mathrm{ens})}\right)
=
\frac{1 + \langle \cos\theta \rangle^2 + \langle \sin\theta \rangle^2}{2}.
\]

This shows that the loss of purity is directly governed by the reduction of phase coherence in the ensemble. When $|\mu|=1$, the state remains pure, while for $|\mu|<1$ the ensemble becomes mixed.

For this reason, purity will serve as a fundamental diagnostic tool in the numerical simulator, complementing correlation observables and preparing the ground for more advanced quantities such as fidelity and entanglement measures.

%========================================================
\subsection{Fidelity with Respect to Bell States}

Having introduced reduced density operators and purity as measures of local information and mixedness, we now require a quantitative tool to evaluate how close a given quantum state is to a reference target state. This role is fulfilled by the notion of fidelity.

In general, the fidelity between two quantum states described by density operators $\rho$ and $\sigma$ is defined as
\[
F(\rho,\sigma)
=
\left(\operatorname{Tr}\sqrt{\sqrt{\rho}\,\sigma\,\sqrt{\rho}}\right)^2.
\]

This quantity satisfies
\[
0 \le F(\rho,\sigma) \le 1,
\]
with $F=1$ if and only if $\rho = \sigma$, and $F=0$ for perfectly distinguishable (orthogonal) states. Fidelity therefore provides a measure of similarity between quantum states.

\medskip

In the present context, the reference state is a pure Bell state, and the fidelity simplifies considerably. Let
\[
\sigma = |\phi\rangle\langle\phi|,
\]
with $|\phi\rangle$ a normalized state vector. Then one obtains (see Appendix~B for a detailed derivation)
\[
F(\rho, |\phi\rangle)
=
\langle \phi | \rho | \phi \rangle.
\]

This expression admits a direct physical interpretation: it corresponds to the probability that the state $\rho$ is found in the state $|\phi\rangle$ upon measurement.

\medskip

In this work, we focus on the Bell state
\[
|\Psi^+\rangle = \frac{|01\rangle + |10\rangle}{\sqrt{2}},
\]
and define the fidelity of a generic state $\rho$ with respect to this target as
\[
F_{\Psi^+}(\rho)
=
\langle \Psi^+ | \rho | \Psi^+ \rangle.
\]

\medskip

We now analyze this quantity for the phase-evolved Bell-like state introduced in Section 3.4:
\[
|\psi(\theta)\rangle = \frac{|01\rangle + e^{i\theta}|10\rangle}{\sqrt{2}},
\]
with associated density operator
\[
\rho(\theta) = |\psi(\theta)\rangle\langle\psi(\theta)|.
\]

The fidelity with respect to $|\Psi^+\rangle$ is given by
\[
F_{\Psi^+}(\rho(\theta))
=
|\langle \Psi^+ | \psi(\theta) \rangle|^2.
\]

A direct computation yields
\[
\langle \Psi^+ | \psi(\theta) \rangle
=
\frac{1}{2}(1 + e^{i\theta}),
\]
and therefore
\[
F_{\Psi^+}(\rho(\theta))
=
\frac{1 + \cos\theta}{2}.
\]

\medskip

This result shows that the relative phase $\theta$ modifies the overlap with the reference Bell state, even though the global state remains pure for all values of $\theta$. In particular,
\[
F_{\Psi^+} = 1 \quad \text{for } \theta = 0,
\qquad
F_{\Psi^+} = 0 \quad \text{for } \theta = \pi.
\]

Thus, fidelity captures variations in global coherence that are not reflected in the purity or in the reduced density operators.

\medskip
For ensemble states arising from random phase fluctuations, described by
\[
\rho_{AB}^{(\mathrm{ens})}
=
\frac{1}{N}\sum_{k=1}^{N}
|\psi(\theta_k)\rangle\langle\psi(\theta_k)|,
\]
the fidelity can be computed using the linearity of the expectation value:
\[
F_{\Psi^+}(\rho_{AB}^{(\mathrm{ens})})
=
\frac{1}{N}
\sum_{k=1}^{N}
F_{\Psi^+}(\rho(\theta_k)).
\]

Introducing the ensemble coherence parameter
\[
\mu = \langle e^{i\theta} \rangle,
\]
one finds
\[
F_{\Psi^+}
=
\frac{1}{2}\left(1 + \mathrm{Re}(\mu)\right)
=
\frac{1}{2}\left(1 + \langle \cos\theta \rangle\right).
\]

\medskip

This shows that fidelity directly quantifies the loss of coherence induced by phase fluctuations. In particular, for a uniform distribution over $[0,2\pi]$, one has $\langle \cos\theta \rangle = 0$, leading to
\[
F_{\Psi^+} = \frac{1}{2}.
\]

\medskip

In summary, fidelity provides a quantitative measure of how closely a given state approximates the target Bell state. Unlike purity, which characterizes mixedness, fidelity is sensitive to coherent deformations of the state and therefore plays a central role in the validation and interpretation of the numerical simulations.

\paragraph{Simulation notes.}
Within the simulator architecture, fidelity is introduced as a general-purpose quantitative tool for comparing quantum states in a basis-independent manner. Given a numerically generated state—either in pure vector form or as a density operator—fidelity provides a scalar measure of proximity to a chosen reference state, typically a Bell state in the present framework. 

From an implementation standpoint, this requires supporting both pure-state overlaps and density-matrix evaluations, allowing the same routine to be applied uniformly across different modeling regimes (deterministic evolution, ensemble averaging, or noisy channels). 

Unlike correlation observables, which probe specific measurement directions, fidelity captures a global property of the state and therefore serves as a compact diagnostic for validating transformations, monitoring deviations from ideal entanglement, and benchmarking different physical models within the simulator pipeline.

%========================================================
%========================================================
\subsection{Concurrence as an Entanglement Measure}

Reduced density operators, purity, and fidelity provide complementary information about a bipartite quantum state. Reduced states describe locally accessible information, purity quantifies mixedness, and fidelity measures proximity to a chosen reference state. However, none of these quantities alone provides a direct measure of bipartite entanglement. In particular, fidelity with respect to a Bell state may change under local unitary evolution even when the amount of entanglement remains unchanged. For this reason, a genuine entanglement measure is required.

For two-qubit systems, one of the most widely used entanglement measures is the \emph{concurrence}. It is specifically designed to quantify bipartite entanglement and satisfies the fundamental bounds
\[
0 \le C \le 1,
\]
where
\[
C = 0
\]
corresponds to a separable state, while
\[
C = 1
\]
corresponds to a maximally entangled two-qubit state.

\medskip

We first consider the case of a pure two-qubit state
\[
|\psi\rangle
=
a|00\rangle + b|01\rangle + c|10\rangle + d|11\rangle,
\]
with complex coefficients satisfying
\[
|a|^2 + |b|^2 + |c|^2 + |d|^2 = 1.
\]
For such a state, the concurrence is given by
\[
C(|\psi\rangle) = 2|ad - bc|.
\]

This expression immediately distinguishes separable from entangled states. Indeed, if
\[
|\psi\rangle = |\phi_A\rangle \otimes |\phi_B\rangle,
\]
then the coefficients satisfy
\[
ad - bc = 0,
\]
and therefore
\[
C(|\psi\rangle)=0.
\]
Conversely, for maximally entangled Bell states the concurrence is equal to one.

\medskip

For example, for the Bell state
\[
|\Psi^+\rangle = \frac{|01\rangle + |10\rangle}{\sqrt{2}},
\]
one has
\[
a=0,\qquad b=\frac{1}{\sqrt{2}},\qquad c=\frac{1}{\sqrt{2}},\qquad d=0,
\]
and thus
\[
C(|\Psi^+\rangle)
=
2\left|0 - \frac{1}{2}\right|
=
1.
\]

The same conclusion holds for the phase-evolved Bell-like state
\[
|\psi(\theta)\rangle
=
\frac{|01\rangle + e^{i\theta}|10\rangle}{\sqrt{2}}.
\]
In this case,
\[
a=0,\qquad b=\frac{1}{\sqrt{2}},\qquad c=\frac{e^{i\theta}}{\sqrt{2}},\qquad d=0,
\]
so that
\[
C(|\psi(\theta)\rangle)
=
2\left|-\frac{e^{i\theta}}{2}\right|
=
1.
\]

This is an important result for the present model. The relative phase \(\theta\) changes the overlap with the reference Bell state \( |\Psi^+\rangle \), and therefore modifies the fidelity introduced in Section 4.3. However, it does \emph{not} change the amount of entanglement. The state remains maximally entangled for all values of \(\theta\). This reflects the fact that the phase evolution considered in the reduced model is generated by local unitary transformations, and bipartite entanglement is invariant under local unitaries.

\medskip

For pure bipartite states, concurrence is also directly connected to the reduced density operators introduced in Section 4.1 and to the purity discussed in Section 4.2. In particular, if
\[
\rho_A = \mathrm{Tr}_B(|\psi\rangle\langle\psi|),
\qquad
\rho_B = \mathrm{Tr}_A(|\psi\rangle\langle\psi|),
\]
then one may write
\[
C(|\psi\rangle)
=
2\sqrt{\det(\rho_A)}
=
2\sqrt{\det(\rho_B)},
\]
or equivalently,
\[
C(|\psi\rangle)
=
\sqrt{2\left(1-\gamma(\rho_A)\right)}
=
\sqrt{2\left(1-\gamma(\rho_B)\right)}.
\]

This relation shows that, for pure two-qubit states, entanglement is directly reflected in the mixedness of the reduced states. In particular, if a pure bipartite state is separable, then the reduced states are pure and
\[
\gamma(\rho_A)=\gamma(\rho_B)=1,
\]
which implies
\[
C=0.
\]
By contrast, for a maximally entangled Bell state,
\[
\rho_A=\rho_B=\frac{I}{2},
\]
and therefore
\[
\gamma(\rho_A)=\gamma(\rho_B)=\frac{1}{2},
\]
which yields
\[
C=1.
\]

Thus, Sections 4.1 and 4.2 naturally prepare the ground for concurrence: reduced mixedness of a subsystem, when the global state is pure, is a direct signature of bipartite entanglement.

\medskip

We now turn to the more general case of mixed two-qubit states, described by a density operator
\[
\rho \in \mathbb{C}^{4\times 4}.
\]
For mixed states, concurrence is defined through the \emph{spin-flipped} matrix
\[
\tilde{\rho}
=
(\sigma_y \otimes \sigma_y)\,\rho^*\,(\sigma_y \otimes \sigma_y),
\]
where \(\rho^*\) denotes the complex conjugate of \(\rho\) in the computational basis. One then considers the eigenvalues of the matrix
\[
\rho \tilde{\rho}.
\]
Let \(\lambda_1,\lambda_2,\lambda_3,\lambda_4\) denote the square roots of these eigenvalues, ordered such that
\[
\lambda_1 \ge \lambda_2 \ge \lambda_3 \ge \lambda_4 \ge 0.
\]
The concurrence of \(\rho\) is then defined as
\[
C(\rho)
=
\max\!\left(0,\lambda_1-\lambda_2-\lambda_3-\lambda_4\right).
\]

This formula provides a complete entanglement measure for arbitrary two-qubit states, both pure and mixed. In the pure-state limit, it reduces to the simpler expression
\[
C(|\psi\rangle)=2|ad-bc|.
\]

\medskip

In the present simulator, this mixed-state formulation becomes essential when the global state is no longer described by a single state vector. This occurs, for example, when one introduces ensemble averaging over random phases,
\[
\rho_{AB}^{(\mathrm{ens})}
=
\frac{1}{N}
\sum_{k=1}^{N}
|\psi(\theta_k)\rangle\langle\psi(\theta_k)|,
\]
or more generally when propagation effects lead to decoherence and partial loss of phase information.

For the random-phase ensemble considered in this work, the density operator depends on the coherence parameter
\[
\mu = \langle e^{i\theta}\rangle,
\]
introduced in Sections 4.2 and 4.3. In the computational basis, the ensemble state has the form
\[
\rho_{AB}^{(\mathrm{ens})}
=
\frac{1}{2}|01\rangle\langle 01|
+
\frac{1}{2}|10\rangle\langle 10|
+
\frac{\mu}{2}|10\rangle\langle 01|
+
\frac{\mu^*}{2}|01\rangle\langle 10|.
\]

For this family of states, the concurrence can be computed analytically and is found to be
\[
C\!\left(\rho_{AB}^{(\mathrm{ens})}\right) = |\mu|.
\]

This result is particularly significant because it reveals that the same coherence parameter controlling mixedness and Bell-state fidelity also controls entanglement degradation in the reduced random-phase model. Indeed, from the previous sections we already found that
\[
\gamma\!\left(\rho_{AB}^{(\mathrm{ens})}\right)
=
\frac{1+|\mu|^2}{2},
\qquad
F_{\Psi^+}
=
\frac{1+\operatorname{Re}(\mu)}{2}.
\]

%========================================================
%========================================================
\subsection{Interplay Between Purity, Fidelity, and Concurrence}

The previous sections introduced three quantitative descriptors of bipartite quantum states: purity, fidelity, and concurrence. Each of these quantities captures a different physical aspect of the state. Purity, defined as \(\gamma(\rho)=\operatorname{Tr}(\rho^2)\), characterizes mixedness, fidelity measures similarity to a chosen reference state, and concurrence provides a direct quantification of bipartite entanglement.

In general, these quantities are independent. However, within the random-phase ensemble model considered in this work, they become closely related through a single complex parameter, namely the ensemble coherence
\[
\mu = \langle e^{i\theta} \rangle
=
\langle \cos\theta \rangle + i \langle \sin\theta \rangle.
\]
This parameter encapsulates the effect of phase fluctuations on the global quantum state.

\medskip

For the family of ensemble states introduced in Section 4.4, the relevant quantities
(purity, concurrence, and Bell-state fidelity) are all functions of the same coherence
parameter \(\mu\), but depend on it in qualitatively different ways.

\medskip

Purity depends on the squared magnitude \( |\mu|^2 \), and therefore quantifies the overall degree of coherence retained after ensemble averaging through
\[
\gamma(\rho) = \frac{1 + |\mu|^2}{2}.
\]
Concurrence depends linearly on \( |\mu| \), and directly measures the amount of bipartite entanglement present in the state. By contrast, the fidelity with respect to the Bell state \( |\Psi^+\rangle \) depends only on the real part \( \operatorname{Re}(\mu) \), reflecting the overlap with a specific reference state rather than the total coherence or entanglement.

\medskip

This distinction has important consequences. Since fidelity depends on the projection of \(\mu\) onto the real axis, it does not provide a complete characterization of entanglement. In particular, states with the same concurrence may exhibit different fidelities with respect to \( |\Psi^+\rangle \). This behavior is clearly observed in the deterministic phase model of Experiment 1, where the parameter \(\mu\) lies on the unit circle. In that case, the magnitude \( |\mu| = 1 \) remains constant, and therefore the concurrence is equal to one for all values of \(\theta\), while the fidelity varies as a function of the phase.

\medskip

In contrast, the random-phase ensemble of Experiment 2 leads to a reduction of the magnitude \( |\mu| \), reflecting the loss of phase coherence under statistical averaging. As a consequence, both purity and concurrence decrease, with purity quantified by \(\gamma(\rho)\). In this regime, the reduction of fidelity is associated with a genuine loss of entanglement, rather than a simple change of phase relative to a fixed reference state.

\medskip

A useful geometric interpretation can be obtained by representing \(\mu\) as a point in the complex plane. The concurrence is proportional to the distance of this point from the origin, while the fidelity depends on its projection onto the real axis. Under deterministic phase evolution, \(\mu\) moves along the unit circle, preserving its magnitude but changing its real part. Under ensemble averaging, \(\mu\) contracts toward the origin, leading to a simultaneous reduction of coherence, purity (as measured by \(\gamma(\rho)\)), and entanglement.

\medskip

The limiting cases illustrate these behaviors clearly. When
\[
|\mu| = 1,
\]
the ensemble retains full phase coherence, the global state is pure, and therefore
\[
\gamma(\rho) = 1,
\]
while the concurrence attains its maximal value. When instead
\[
|\mu| = 0,
\]
all off-diagonal coherence terms vanish, the state reduces to a classical mixture, for which
\[
\gamma(\rho) = \frac{1}{2},
\]
and the concurrence becomes zero. Intermediate values \(0 < |\mu| < 1\) correspond to partially coherent ensembles, characterized by mixed states with reduced but non-vanishing entanglement.

\medskip

It is important to emphasize that, while purity, fidelity, and concurrence are general concepts in quantum information theory, the simple relations above arise from the specific structure of the random-phase ensemble considered here. In more general situations, these quantities are not necessarily determined by a single parameter.

\medskip

In summary, the coherence parameter \(\mu\) provides a unifying description of the ensemble states studied in this work. Its magnitude controls both mixedness and entanglement, while its phase determines the observable correlations and the overlap with a chosen reference state. This interplay clarifies the distinct roles of purity, fidelity, and concurrence, and highlights the necessity of using a genuine entanglement measure to fully characterize bipartite quantum states.

%========================================================
%========================================================
%========================================================
\section{Numerical Experiments and Results}

\subsection{Experiment 1: Deterministic Phase Sweep}

\subsubsection{Objective}

The first experiment investigates the effect of a controlled relative phase \(\theta\) on a maximally entangled two-qubit state. The objective is to characterize how both correlation observables and state-based quantities evolve under the reduced unitary fiber model.

\subsubsection{Model}

The initial Bell state is chosen as
\[
|\Psi^{+}\rangle
=
\frac{1}{\sqrt{2}}
\left(
|01\rangle + |10\rangle
\right).
\]

After propagation through the effective phase model, the state becomes
\[
|\psi(\theta)\rangle
=
\frac{1}{\sqrt{2}}
\left(
|01\rangle + e^{i\theta}|10\rangle
\right).
\]

The correlation observables evaluated in the simulator are
\[
c_{xx} = \langle \sigma_x \otimes \sigma_x \rangle,\qquad
c_{yy} = \langle \sigma_y \otimes \sigma_y \rangle,\qquad
c_{zz} = \langle \sigma_z \otimes \sigma_z \rangle.
\]

\subsubsection{Analytical Predictions}

For the state \( |\psi(\theta)\rangle \), direct evaluation gives
\[
c_{xx} = \cos\theta,\qquad
c_{yy} = \cos\theta,\qquad
c_{zz} = -1.
\]

Since the evolution is generated by local unitary transformations, the global state remains pure, i.e.
\[
\gamma(\rho(\theta)) = \operatorname{Tr}(\rho(\theta)^2)=1.
\]

The reduced density operators remain maximally mixed:
\[
\rho_A(\theta)=\rho_B(\theta)=\frac{I}{2}, \qquad
\gamma(\rho_A)=\gamma(\rho_B)=\operatorname{Tr}(\rho_A^2)=\operatorname{Tr}(\rho_B^2)=\frac{1}{2}.
\]

The fidelity with respect to the reference Bell state \( |\Psi^+\rangle \) is
\[
F_{\Psi^+}(\rho(\theta))
=
\frac{1+\cos\theta}{2}.
\]

Finally, the concurrence of the state is constant and equal to
\[
C(\rho(\theta)) = 1,
\]
reflecting the fact that bipartite entanglement is invariant under local unitary transformations.

\subsubsection{Numerical Results}

The numerical experiment performs a sweep over \(\theta \in [0,2\pi]\), computing both correlation observables and state-characterization quantities.

\begin{figure}[h]
    \centering
    \includegraphics[width=0.75\linewidth]{../Images/Experiment_1/experiment_1_correlations_validation.png}
    \caption{
    Comparison between numerical simulation and analytical predictions for the correlation observables 
    \(c_{xx}\), \(c_{yy}\), and \(c_{zz}\) as functions of the relative phase \(\theta\). 
    The agreement confirms the correctness of the reduced phase model implementation.
    }
    \label{fig:experiment1_validation}
\end{figure}

\begin{figure}[h]
    \centering
    \includegraphics[width=0.75\linewidth]{../Images/Experiment_1/experiment_1_state_metrics.png}
    \caption{
    Numerical and analytical comparison of state-characterization quantities for the phase-evolved state.
    The figure shows global purity \(\gamma(\rho_{AB})\), reduced purities \(\gamma(\rho_A)\) and \(\gamma(\rho_B)\), fidelity with respect to \( |\Psi^+\rangle \), and concurrence.
    While fidelity varies with \(\theta\), concurrence remains constant and equal to one.
    }
    \label{fig:experiment1_state_metrics}
\end{figure}

\subsubsection{Discussion}

The numerical results are in full agreement with the analytical predictions:
\begin{itemize}
    \item \(c_{xx}\) and \(c_{yy}\) follow the expected cosine dependence,
    \item \(c_{zz}\) remains constant at \(-1\),
    \item the global purity remains equal to \(\gamma(\rho_{AB})=1\),
    \item the reduced purities remain equal to \(\gamma(\rho_A)=\gamma(\rho_B)=1/2\),
    \item the fidelity follows \(F_{\Psi^+}=(1+\cos\theta)/2\),
    \item the concurrence remains constant at \(C=1\).
\end{itemize}

These results confirm that the fiber acts as an effective relative phase shifter. While the global state remains pure and maximally entangled, both the observable correlations and the fidelity with respect to a fixed Bell reference vary continuously with the phase.

A key observation is that fidelity and concurrence exhibit fundamentally different behaviors. The fidelity depends explicitly on the phase \(\theta\), reflecting the overlap with a fixed reference state, whereas the concurrence remains invariant. This highlights that fidelity does not provide a direct measure of entanglement.

Therefore, although the phase evolution modifies observable correlations and the apparent similarity to a chosen Bell state, it does not alter the underlying entanglement structure. This demonstrates that purity (quantified by \(\gamma(\rho)\)) and fidelity alone are insufficient to characterize bipartite entanglement, and motivates the introduction of concurrence as a genuine entanglement measure.

\subsubsection{Physical Interpretation}

The results obtained in this experiment admit a direct physical interpretation in terms of fiber-induced birefringence. In optical fibers, the horizontal and vertical polarization components generally propagate with different effective refractive indices, leading to a relative phase accumulation.

After compensation of global polarization rotations, the dominant residual effect is a differential phase shift between the two polarization modes. In the quantum description, this corresponds to a local unitary transformation equivalent to a rotation around the Z-axis of the Bloch sphere.

As a consequence, the initial Bell state evolves into a phase-dependent superposition while preserving its entangled nature. The invariance of the concurrence reflects the fact that no information is lost during propagation, and the system remains in a pure state.

At the same time, the phase dependence of the transverse correlations \(c_{xx}\) and \(c_{yy}\) demonstrates that measurable quantities are sensitive to this relative phase. This highlights the operational relevance of the phase parameter \(\theta\), which becomes the key variable controlling observable correlations in the absence of noise.

This experiment therefore establishes the minimal physical model of fiber propagation: a deterministic unitary transformation that modifies observable correlations without degrading entanglement.
%========================================================
%========================================================
\subsection{Experiment 2: Random Phase Ensemble}

\subsubsection{Objective}

The second experiment extends the deterministic phase model by introducing random phase realizations \(\theta_k\), modeling stochastic birefringence effects in realistic optical fibers.

Instead of a single pure state, the system is described by an ensemble
\[
\{ |\psi(\theta_k)\rangle \}_{k=1}^{N},
\]
and by the corresponding averaged density operator.

The objective is to characterize how ensemble averaging modifies global coherence, correlation observables, purity, fidelity, and entanglement.

\subsubsection{Ensemble Model}

Each realization is given by
\[
|\psi(\theta_k)\rangle
=
\frac{1}{\sqrt{2}}
\left(
|01\rangle + e^{i\theta_k}|10\rangle
\right),
\]
and the ensemble state is
\[
\rho_{\mathrm{ensemble}}
=
\frac{1}{N}
\sum_{k=1}^{N}
|\psi(\theta_k)\rangle\langle \psi(\theta_k)|.
\]

\subsubsection{Analytical Predictions}

For the ensemble state, it is convenient to introduce the ensemble coherence parameter
\[
\mu = \langle e^{i\theta} \rangle
=
\langle \cos\theta \rangle + i \langle \sin\theta \rangle.
\]

The ensemble-averaged correlations satisfy
\[
c_{xx} = \mathrm{Re}(\mu) = \langle \cos\theta \rangle,\qquad
c_{yy} = \mathrm{Re}(\mu) = \langle \cos\theta \rangle,\qquad
c_{zz} = -1.
\]

The global purity is
\[
\gamma(\rho_{\mathrm{ensemble}})
=
\operatorname{Tr}(\rho_{\mathrm{ensemble}}^2)
=
\frac{1 + |\mu|^2}{2}
=
\frac{1 + \langle \cos\theta \rangle^2 + \langle \sin\theta \rangle^2}{2}.
\]

The reduced states remain maximally mixed:
\[
\rho_A = \rho_B = \frac{I}{2},
\]
and therefore
\[
\gamma(\rho_A)=\gamma(\rho_B)
=
\operatorname{Tr}(\rho_A^2)=\operatorname{Tr}(\rho_B^2)=\frac{1}{2}.
\]

The fidelity with respect to \( |\Psi^+\rangle \) is
\[
F_{\Psi^+}
=
\frac{1 + \mathrm{Re}(\mu)}{2}
=
\frac{1 + \langle \cos\theta \rangle}{2}.
\]

Finally, the concurrence of the ensemble state is
\[
C(\rho_{\mathrm{ensemble}}) = |\mu|,
\]
showing that entanglement is directly controlled by the magnitude of the ensemble coherence parameter.

\subsubsection{Numerical Results}

Three representative phase distributions are considered:
\begin{itemize}
    \item constant phase,
    \item uniform distribution,
    \item Gaussian distribution.
\end{itemize}

Each case is summarized in a single figure combining:
\begin{itemize}
    \item sampled phase statistics,
    \item pure-state correlation trend,
    \item ensemble correlations,
    \item purity, fidelity, and concurrence metrics.
\end{itemize}

\begin{figure}[h]
    \centering
    \includegraphics[width=0.9\linewidth]{../Images/Experiment_2/experiment_2_constant_distribution.png}
    \caption{
    Summary of Experiment 2 for a constant phase distribution. The ensemble reduces to a single deterministic realization, and concurrence remains equal to one.
    }
    \label{fig:experiment2_constant}
\end{figure}

\begin{figure}[h]
    \centering
    \includegraphics[width=0.9\linewidth]{../Images/Experiment_2/experiment_2_uniform_distribution.png}
    \caption{
    Summary of Experiment 2 for a uniform phase distribution. Phase averaging suppresses transverse correlations, reduces fidelity, and drives concurrence to zero.
    }
    \label{fig:experiment2_uniform}
\end{figure}

\begin{figure}[h]
    \centering
    \includegraphics[width=0.9\linewidth]{../Images/Experiment_2/experiment_2_gaussian_distribution.png}
    \caption{
    Summary of Experiment 2 for a Gaussian phase distribution. Residual coherence depends on the phase variance, leading to intermediate values of concurrence.
    }
    \label{fig:experiment2_gaussian}
\end{figure}

\subsubsection{Discussion}

The results show a transition from coherent deterministic behavior to mixed-state ensemble behavior:
\begin{itemize}
    \item for constant \(\theta\), the system reproduces Experiment 1,
    \item for distributed \(\theta\), transverse correlations are reduced,
    \item for broad phase distributions, \( |\langle e^{i\theta}\rangle| \to 0 \), leading to vanishing ensemble coherence.
\end{itemize}

At the same time:
\begin{itemize}
    \item the global purity decreases from \(\gamma(\rho_{\mathrm{ensemble}})=1\),
    \item the reduced purities remain fixed at \(\gamma(\rho_A)=\gamma(\rho_B)=1/2\),
    \item the fidelity with respect to \( |\Psi^+\rangle \) decreases,
    \item the concurrence decreases from \(1\) to \(0\).
\end{itemize}

This behavior highlights a fundamental difference with respect to Experiment 1. In the deterministic case, variations in fidelity were not associated with changes in entanglement. Here, instead, the decrease of fidelity is accompanied by a genuine loss of entanglement, as quantified by concurrence.

\subsubsection{Physical Interpretation}

Each realization remains a pure maximally entangled state. The mixedness of the ensemble arises exclusively from statistical averaging over unresolved phase fluctuations.

The coherence parameter \(\mu = \langle e^{i\theta} \rangle\) plays a central role: its real part determines observable correlations and fidelity, while its magnitude controls the amount of entanglement.

This demonstrates the transition from a pure-state description to an effective mixed-state description, where coherence is progressively lost at the observable level and entanglement is gradually degraded.

In contrast to the deterministic case, where local unitary transformations preserve entanglement, the ensemble averaging process leads to a genuine reduction of entanglement, reflecting the loss of phase information in realistic propagation scenarios.

%========================================================
%========================================================
\subsection{Experiment 3: Depolarizing Channel}

\subsubsection{Objective}

The purpose of this experiment is to study the degradation of a maximally entangled Bell state under the action of an effective two-qubit depolarizing channel. Unlike the previous experiments, where the dynamics was generated by deterministic or ensemble-averaged unitary phase evolution, the present model introduces a genuinely non-unitary transformation at the level of the density operator.

The main objectives are:
\begin{itemize}
    \item to validate the implementation of the depolarizing channel introduced in Section 3.8;
    \item to quantify the effect of depolarization on global purity, reduced purities, fidelity with respect to \(|\Psi^+\rangle\), concurrence, and Pauli correlations;
    \item to compare numerical results against the corresponding analytical predictions over the full range \(p \in [0,1]\).
\end{itemize}

This experiment therefore represents the first explicit implementation of effective noisy quantum evolution in the simulator.

\subsubsection{Model}

The reference input state is the Bell state
\[
|\Psi^+\rangle = \frac{|01\rangle + |10\rangle}{\sqrt{2}},
\]
with associated density operator
\[
\rho_0 = |\Psi^+\rangle\langle\Psi^+|.
\]

The depolarizing model is defined by the map
\[
\rho(p) = (1-p)\rho_0 + \frac{p}{4}I_4,
\qquad p \in [0,1],
\]
where \(I_4\) is the identity operator on the two-qubit Hilbert space. The parameter \(p\) quantifies the strength of the depolarizing contribution.

This model interpolates between two limiting regimes:
\[
\rho(0)=\rho_0,
\qquad
\rho(1)=\frac{I_4}{4}.
\]
Thus, for \(p=0\) the output state is the initial Bell state, while for \(p=1\) the output is the maximally mixed two-qubit state.

At the MATLAB level, the experiment is implemented by sweeping a set of values \(p_1,\dots,p_N\), applying the function \texttt{depolarizing\_channel\_two\_qubits} to \(\rho_0\), and evaluating for each output state:
\begin{itemize}
    \item the global purity \(\gamma_{AB} = \operatorname{Tr}(\rho^2)\),
    \item the reduced purities \(\gamma_A\) and \(\gamma_B\),
    \item the fidelity \(F_{\Psi^+}\) with respect to the reference Bell state,
    \item the concurrence \(C\),
    \item the Pauli correlations \(c_{xx}\), \(c_{yy}\), and \(c_{zz}\).
\end{itemize}

\subsubsection{Analytical Predictions}

For the family of states
\[
\rho(p) = (1-p)|\Psi^+\rangle\langle\Psi^+| + \frac{p}{4}I_4,
\]
the analytical predictions used for validation are:
\[
\gamma_{AB}(p)=1-\frac{3}{2}p+\frac{3}{4}p^2,
\]
\[
\gamma_A(p)=\gamma_B(p)=\frac{1}{2},
\]
\[
F_{\Psi^+}(p)=1-\frac{3}{4}p,
\]
\[
C(p)=\max\left(0,\,1-\frac{3}{2}p\right),
\]
and
\[
c_{xx}(p)=1-p,\qquad
c_{yy}(p)=1-p,\qquad
c_{zz}(p)=-(1-p).
\]

These expressions show that the depolarizing channel affects the different state descriptors in qualitatively distinct ways. The global purity decreases quadratically, the fidelity decreases linearly, and the concurrence exhibits threshold behavior, vanishing at
\[
p=\frac{2}{3}.
\]
By contrast, the reduced purities remain constant at \(1/2\), reflecting the fact that the reduced states remain maximally mixed throughout the entire family.

\subsubsection{Numerical Results}

The experiment was executed by sampling the depolarization parameter uniformly over the interval \([0,1]\). For each sampled value, the simulator constructed the output density matrix, computed all target observables, and compared the numerical values with the analytical ones.

The console-monitoring summary reported maximum absolute errors at machine-precision level for all tested quantities, confirming exact numerical agreement with theory up to floating-point tolerance. In particular, the comparison showed negligible discrepancies for:
\begin{itemize}
    \item global purity \(\gamma_{AB}\),
    \item reduced purities \(\gamma_A\) and \(\gamma_B\),
    \item fidelity \(F_{\Psi^+}\),
    \item concurrence \(C\),
    \item correlations \(c_{xx}\), \(c_{yy}\), and \(c_{zz}\).
\end{itemize}

\begin{figure}[h]
    \centering
    \includegraphics[width=\textwidth]{../Images/Experiment_3/experiment_3_depolarizing_channel_summary.png}
    \caption{Experiment 3: numerical and analytical behavior of state metrics and Pauli correlations under the two-qubit depolarizing channel applied to the Bell state \(|\Psi^+\rangle\). The agreement is exact up to numerical precision.}
    \label{fig:experiment3_depolarizing_summary}
\end{figure}

\subsubsection{Discussion}

This experiment validates the first non-unitary extension of the simulator. In the previous experiments, mixed states arose from averaging over unresolved unitary realizations. Here, instead, the state is evolved directly through an effective quantum channel.

The results confirm several physically important points. First, depolarization reduces global coherence and entanglement while preserving the maximally mixed character of the local reduced states. Second, concurrence and purity do not encode the same information: while the global purity decreases smoothly over the entire interval, concurrence vanishes already at \(p=2/3\), after which the state remains mixed but separable. Third, the correlation observables decay linearly in magnitude, providing a direct operational signature of the channel action.

From the point of view of simulator development, Experiment 3 is a crucial milestone because it verifies that the code correctly handles channel-based evolution, mixed-state metrics, and analytical benchmarking within a unified framework.

\subsubsection{Physical Interpretation}

The depolarizing channel introduced here is not intended as a complete microscopic fiber model. Rather, it serves as a first effective description of coherence loss when polarization becomes partially inaccessible due to unresolved physical mechanisms such as frequency dependence, temporal walk-off, or polarization-mode dispersion.

In this sense, the channel provides a compact phenomenological model of entanglement degradation. Its action drives the initially pure Bell state toward the maximally mixed state, while progressively suppressing both nonlocal correlations and entanglement.

This makes the experiment especially valuable as a bridge between the abstract theory of quantum channels and the future development of more physically detailed fiber models. In particular, it establishes a baseline against which more refined descriptions---for example segmented random-unitary models or frequency-dependent propagation models---can later be compared.

%========================================================
%========================================================
%========================================================
\appendix
\section{Explicit Derivation of Reduced Density Operators and Comparison of Global States}

This appendix provides an explicit derivation of the reduced density operators for the phase-evolved Bell-like state and illustrates, through a concrete comparison, why identical local reduced states may correspond to different global physical situations.

We consider the two-qubit state
\[
|\psi(\theta)\rangle = \frac{|01\rangle + e^{i\theta}|10\rangle}{\sqrt{2}},
\]
written in the computational basis
\[
\{|00\rangle, |01\rangle, |10\rangle, |11\rangle\}.
\]

Its vector representation is
\[
|\psi(\theta)\rangle =
\frac{1}{\sqrt{2}}
\begin{pmatrix}
0 \\
1 \\
e^{i\theta} \\
0
\end{pmatrix}.
\]

The associated density operator is
\[
\rho_{AB}(\theta) = |\psi(\theta)\rangle\langle\psi(\theta)|,
\]
which expands as
\[
\rho_{AB}(\theta)
=
\frac{1}{2}
\Big(
|01\rangle\langle 01|
+
e^{-i\theta}|01\rangle\langle 10|
+
e^{i\theta}|10\rangle\langle 01|
+
|10\rangle\langle 10|
\Big).
\]

In matrix form,
\[
\rho_{AB}(\theta)
=
\frac{1}{2}
\begin{pmatrix}
0 & 0 & 0 & 0 \\
0 & 1 & e^{-i\theta} & 0 \\
0 & e^{i\theta} & 1 & 0 \\
0 & 0 & 0 & 0
\end{pmatrix}.
\]

We now compute the reduced density operator of subsystem $A$ by taking the partial trace over subsystem $B$:
\[
\rho_A(\theta) = \mathrm{Tr}_B(\rho_{AB}(\theta)).
\]

Using the explicit formula
\[
\rho_A(\theta)
=
\sum_{b=0}^{1}
(I_A \otimes \langle b|)\,\rho_{AB}(\theta)\,(I_A \otimes |b\rangle),
\]
we evaluate the contribution of each term in $\rho_{AB}(\theta)$ separately.

For the diagonal term $|01\rangle\langle 01|$, with
\[
|01\rangle = |0\rangle_A \otimes |1\rangle_B,
\]
one obtains
\[
\mathrm{Tr}_B(|01\rangle\langle 01|) = |0\rangle\langle 0|.
\]

For the coherence term
\[
|01\rangle\langle 10|
=
(|0\rangle\langle 1|)\otimes(|1\rangle\langle 0|),
\]
one finds
\[
\mathrm{Tr}_B(|01\rangle\langle 10|)
=
|0\rangle\langle 1|\,\mathrm{Tr}(|1\rangle\langle 0|)
=
0,
\]
since
\[
\mathrm{Tr}(|1\rangle\langle 0|)=\langle 0|1\rangle = 0.
\]
The same argument gives
\[
\mathrm{Tr}_B(|10\rangle\langle 01|) = 0.
\]

Finally, for the second diagonal term $|10\rangle\langle 10|$, with
\[
|10\rangle = |1\rangle_A \otimes |0\rangle_B,
\]
one obtains
\[
\mathrm{Tr}_B(|10\rangle\langle 10|) = |1\rangle\langle 1|.
\]

Summing all contributions, the reduced density operator of subsystem $A$ is
\[
\rho_A(\theta)
=
\frac{1}{2}|0\rangle\langle 0|
+
\frac{1}{2}|1\rangle\langle 1|
=
\frac{I}{2}.
\]

By symmetry, the same result holds for subsystem $B$:
\[
\rho_B(\theta) = \frac{I}{2}.
\]

Therefore, although the global density operator depends explicitly on the phase $\theta$, the reduced density operators do not. The phase appears only in off-diagonal coherence terms, and these are eliminated by the partial trace because of the orthogonality relations in the traced subsystem.

\medskip

Setting $\theta = 0$ recovers the Bell state
\[
|\Psi^+\rangle = \frac{|01\rangle + |10\rangle}{\sqrt{2}},
\]
which gives the same reduced density operators
\[
\rho_A = \rho_B = \frac{I}{2}.
\]

This shows explicitly that each subsystem is locally maximally mixed even though the global state is pure.

\medskip

To understand why this does not imply classical ignorance, consider the single-qubit mixed state
\[
\rho^{(\mathrm{class})}
=
\frac{1}{2}|0\rangle\langle 0|
+
\frac{1}{2}|1\rangle\langle 1|
=
\frac{I}{2}.
\]

At the level of subsystem $A$, the reduced state obtained from the entangled two-qubit state coincides with this classical mixture:
\[
\rho_A^{(\mathrm{ent})} = \rho^{(\mathrm{class})}.
\]

Consequently, for any local observable $O_A$,
\[
\mathrm{Tr}(\rho^{(\mathrm{class})} O_A)
=
\mathrm{Tr}(\rho_A^{(\mathrm{ent})} O_A).
\]

Thus, no measurement acting only on subsystem $A$ can distinguish whether the local mixedness arises from a classical mixture or from entanglement with subsystem $B$.

\medskip

At the global level, however, the distinction becomes manifest. Consider the separable mixed state
\[
\rho_{AB}^{(\mathrm{sep})}
=
\frac{1}{2}|01\rangle\langle 01|
+
\frac{1}{2}|10\rangle\langle 10|.
\]

This state has the same reduced density operators as the entangled state,
\[
\mathrm{Tr}_B(\rho_{AB}^{(\mathrm{sep})}) = \frac{I}{2},
\qquad
\mathrm{Tr}_A(\rho_{AB}^{(\mathrm{sep})}) = \frac{I}{2},
\]
but it does not contain the off-diagonal coherence terms
\[
e^{-i\theta}|01\rangle\langle 10|,
\qquad
e^{i\theta}|10\rangle\langle 01|.
\]

One immediate indicator of this difference is the purity
\[
\gamma(\rho) = \operatorname{Tr}(\rho^2).
\]

For the entangled pure state,
\[
\gamma(\rho_{AB}(\theta)) = \operatorname{Tr}(\rho_{AB}(\theta)^2) = 1,
\]
while for the separable mixed state,
\[
\gamma(\rho_{AB}^{(\mathrm{sep})}) = \operatorname{Tr}\big((\rho_{AB}^{(\mathrm{sep})})^2\big) = \frac{1}{2}.
\]

A second distinction appears in the bipartite correlation observables. For the phase-evolved state,
\[
\langle \sigma_x \otimes \sigma_x \rangle = \cos\theta, \qquad
\langle \sigma_y \otimes \sigma_y \rangle = \cos\theta, \qquad
\langle \sigma_z \otimes \sigma_z \rangle = -1,
\]
whereas for the separable mixed state one obtains
\[
\langle \sigma_x \otimes \sigma_x \rangle = 0, \qquad
\langle \sigma_y \otimes \sigma_y \rangle = 0, \qquad
\langle \sigma_z \otimes \sigma_z \rangle = -1.
\]

Thus, the two states are identical from the local point of view but differ at the global level in both purity and correlation structure. This confirms explicitly that reduced density operators encode all locally accessible information, but do not in general capture global quantum coherence or bipartite entanglement.

%==============================================================
\section{Derivation of Fidelity for a Pure Reference State}

We derive the simplified expression for the fidelity when one of the states is pure, and justify the algebraic properties used in the reduction.

\medskip

Let $\rho$ be a density operator on a finite-dimensional Hilbert space, and let
\[
\sigma = |\psi\rangle\langle\psi|
\]
be a pure-state density operator, with $|\psi\rangle$ normalized:
\[
\langle \psi | \psi \rangle = 1.
\]

By construction, $\sigma$ is Hermitian,
\[
\sigma^\dagger = \sigma,
\]
positive semidefinite,
\[
\langle \chi | \sigma | \chi \rangle \ge 0
\quad \text{for all } |\chi\rangle,
\]
and satisfies
\[
\operatorname{Tr}(\sigma) = 1.
\]
Moreover, it is a rank-one operator, since it projects onto the one-dimensional subspace spanned by $|\psi\rangle$.

\medskip

The general definition of fidelity is
\[
F(\rho,\sigma)
=
\left(\operatorname{Tr}\sqrt{\sqrt{\rho}\,\sigma\,\sqrt{\rho}}\right)^2.
\]

Substituting $\sigma = |\psi\rangle\langle\psi|$, one obtains
\[
\sqrt{\rho}\,\sigma\,\sqrt{\rho}
=
\sqrt{\rho}\,|\psi\rangle\langle\psi|\,\sqrt{\rho}.
\]

Define the vector
\[
|\phi\rangle = \sqrt{\rho}\,|\psi\rangle.
\]
Then the operator becomes
\[
\sqrt{\rho}\,\sigma\,\sqrt{\rho}
=
|\phi\rangle\langle\phi|.
\]

This is again a rank-one operator, provided $|\phi\rangle \neq 0$.

\medskip

We now compute the square root of this operator. Since
\[
|\phi\rangle\langle\phi|
\]
is positive semidefinite and rank one, its only nonzero eigenvalue is
\[
\lambda = \langle \phi | \phi \rangle.
\]

Thus, its spectral decomposition is
\[
|\phi\rangle\langle\phi|
=
\lambda\,|\hat{\phi}\rangle\langle\hat{\phi}|,
\]
where
\[
|\hat{\phi}\rangle = \frac{|\phi\rangle}{\|\phi\|}.
\]

The square root is then
\[
\sqrt{|\phi\rangle\langle\phi|}
=
\sqrt{\lambda}\,|\hat{\phi}\rangle\langle\hat{\phi}|.
\]

Taking the trace,
\[
\operatorname{Tr}\sqrt{|\phi\rangle\langle\phi|}
=
\sqrt{\lambda}.
\]

Therefore,
\[
F(\rho,\sigma)
=
\lambda
=
\langle \phi | \phi \rangle.
\]

Recalling the definition of $|\phi\rangle$,
\[
\lambda
=
\langle \psi | \rho | \psi \rangle.
\]

\medskip

We conclude that when $\sigma = |\psi\rangle\langle\psi|$ is a pure state, the fidelity reduces to
\[
F(\rho,|\psi\rangle\langle\psi|)
=
\langle \psi | \rho | \psi \rangle.
\]

\medskip

This result justifies the simplified expression used in Section 4.3. It also shows that the fidelity corresponds to the expectation value of the projector onto $|\psi\rangle$, providing a direct probabilistic interpretation as the probability of obtaining the outcome $|\psi\rangle$ in a projective measurement.
%========================================================
%========================================================
\section{Derivations for Concurrence and Reduced Density Operators}

This appendix collects the main algebraic derivations leading to the compact expressions for concurrence and its relation to reduced density operators.

Let
\[
|\psi\rangle
=
a|00\rangle + b|01\rangle + c|10\rangle + d|11\rangle
\]
be a normalized pure two-qubit state.

It is convenient to introduce the associated coefficient matrix
\[
A_\psi =
\begin{pmatrix}
a & b\\
c & d
\end{pmatrix},
\]
such that
\[
|\psi\rangle
=
\sum_{i,j=0}^{1} (A_\psi)_{ij}\,|i\rangle \otimes |j\rangle.
\]

This representation allows one to rewrite bipartite operations in matrix form.

The global density operator is
\[
\rho_{AB} = |\psi\rangle\langle\psi|.
\]

Expanding in the computational basis,
\[
\rho_{AB}
=
\sum_{i,j,k,l}
(A_\psi)_{ij}(A_\psi)^*_{kl}
\;
|i\rangle\langle k|
\otimes
|j\rangle\langle l|.
\]

The reduced density operator on subsystem \(A\) is
\[
\rho_A = \operatorname{Tr}_B(\rho_{AB})
=
\sum_v
\bigl(I_A \otimes \langle v|\bigr)
\rho_{AB}
\bigl(I_A \otimes |v\rangle\bigr),
\]
where \(\{|v\rangle\}\) is an orthonormal basis of \(\mathcal{H}_B\).

Applying this to the expansion above,
\[
\rho_A
=
\sum_{i,j,k,l}
(A_\psi)_{ij}(A_\psi)^*_{kl}
\left(
\sum_v \langle v|j\rangle \langle l|v\rangle
\right)
|i\rangle\langle k|.
\]

Using the completeness relation
\[
\sum_v |v\rangle\langle v| = I,
\]
one obtains
\[
\sum_v \langle v|j\rangle \langle l|v\rangle = \delta_{jl}.
\]

Therefore,
\[
\rho_A
=
\sum_{i,k,j}
(A_\psi)_{ij}(A_\psi)^*_{kj}
\,|i\rangle\langle k|.
\]

In matrix form, this corresponds to
\[
\rho_A = A_\psi A_\psi^\dagger.
\]

Similarly,
\[
\rho_B = \operatorname{Tr}_A(\rho_{AB})
=
\sum_{j,\ell,i}
(A_\psi)_{ij}(A_\psi)^*_{i\ell}
\,|j\rangle\langle \ell|,
\]
which gives
\[
\rho_B = A_\psi^\dagger A_\psi.
\]


The concurrence of a pure two-qubit state is defined as
\[
C(|\psi\rangle)
=
\left|
\langle \psi | (\sigma_y \otimes \sigma_y)\,|\psi^*\rangle
\right|.
\]

Let
\[
|\psi\rangle =
\begin{pmatrix}
a\\
b\\
c\\
d
\end{pmatrix},
\qquad
|\psi^*\rangle =
\begin{pmatrix}
a^*\\
b^*\\
c^*\\
d^*
\end{pmatrix}.
\]

Using the matrix representation
\[
\sigma_y \otimes \sigma_y =
\begin{pmatrix}
0 & 0 & 0 & -1\\
0 & 0 & 1 & 0\\
0 & 1 & 0 & 0\\
-1 & 0 & 0 & 0
\end{pmatrix},
\]
one finds
\[
(\sigma_y \otimes \sigma_y)|\psi^*\rangle
=
\begin{pmatrix}
- d^*\\
c^*\\
b^*\\
- a^*
\end{pmatrix}.
\]

Taking the inner product with \(\langle \psi|\),
\[
\langle \psi | (\sigma_y \otimes \sigma_y)|\psi^*\rangle
=
- a^* d^* + b^* c^* + c^* b^* - d^* a^*.
\]

Using commutativity of scalar multiplication,
\[
=
2(b^*c^* - a^*d^*)
=
2(ad - bc)^*.
\]

Taking the modulus,
\[
C(|\psi\rangle)=2|ad-bc|.
\]

Introducing the coefficient matrix \(A_\psi\),
\[
\det(A_\psi) = ad - bc,
\]
so that
\[
C(|\psi\rangle)=2|\det(A_\psi)|.
\]


Using
\[
\rho_A = A_\psi A_\psi^\dagger,
\]
one obtains
\[
\det(\rho_A)
=
\det(A_\psi)\det(A_\psi^\dagger).
\]

Since
\[
\det(A_\psi^\dagger) = (\det A_\psi)^*,
\]
it follows that
\[
\det(\rho_A)=|\det(A_\psi)|^2.
\]

Therefore,
\[
C(|\psi\rangle)=2\sqrt{\det(\rho_A)}.
\]

An identical relation holds for \(\rho_B\).


Since $\rho$ is a density operator, it is positive semidefinite and normalized,
\[
\operatorname{Tr}(\rho) = 1.
\]
For any $2\times2$ Hermitian matrix with unit trace, the determinant can be expressed in terms of the purity
\[
\gamma(\rho) = \operatorname{Tr}(\rho^2)
\]
as
\[
\det(\rho)
=
\frac{1 - \gamma(\rho)}{2}.
\]

Substituting into the previous result,
\[
C(|\psi\rangle)
=
\sqrt{2\left(1 - \gamma(\rho_A)\right)}.
\]

This shows explicitly that, for pure bipartite states, entanglement is directly encoded in the mixedness of the reduced subsystems.
%========================================================
%========================================================
\section{Analytical Derivations for Experiment 3: Depolarizing Channel}

This appendix derives the analytical expressions used in Experiment 3 for the family of two-qubit states obtained by applying the depolarizing channel to the Bell state \(|\Psi^+\rangle\).

We consider the reference Bell state
\[
|\Psi^+\rangle = \frac{|01\rangle + |10\rangle}{\sqrt{2}},
\]
with density operator
\[
\rho_0 = |\Psi^+\rangle\langle\Psi^+|.
\]

The depolarized family is defined as
\[
\rho(p) = (1-p)\rho_0 + \frac{p}{4}I_4,
\qquad p \in [0,1].
\]

\subsection*{Density matrix form}

In the computational basis
\[
\{|00\rangle, |01\rangle, |10\rangle, |11\rangle\},
\]
the projector onto \(|\Psi^+\rangle\) is
\[
\rho_0
=
\begin{pmatrix}
0 & 0 & 0 & 0\\
0 & \frac12 & \frac12 & 0\\
0 & \frac12 & \frac12 & 0\\
0 & 0 & 0 & 0
\end{pmatrix}.
\]

Therefore,
\[
\rho(p)
=
(1-p)
\begin{pmatrix}
0 & 0 & 0 & 0\\
0 & \frac12 & \frac12 & 0\\
0 & \frac12 & \frac12 & 0\\
0 & 0 & 0 & 0
\end{pmatrix}
+
\frac{p}{4}
\begin{pmatrix}
1 & 0 & 0 & 0\\
0 & 1 & 0 & 0\\
0 & 0 & 1 & 0\\
0 & 0 & 0 & 1
\end{pmatrix},
\]
that is,
\[
\rho(p)
=
\begin{pmatrix}
\frac p4 & 0 & 0 & 0\\
0 & \frac12-\frac p4 & \frac{1-p}{2} & 0\\
0 & \frac{1-p}{2} & \frac12-\frac p4 & 0\\
0 & 0 & 0 & \frac p4
\end{pmatrix}.
\]

\subsection*{Global purity}

The global purity is
\[
\gamma_{AB}(p) = \operatorname{Tr}(\rho(p)^2).
\]

A convenient way to compute it is through the eigenvalues of \(\rho(p)\). Since \(\rho_0\) is a rank-one projector and \(I_4/4\) is proportional to the identity, the Bell-state direction remains an eigenvector of \(\rho(p)\), with eigenvalue
\[
\lambda_1 = (1-p)\cdot 1 + \frac p4 = 1-\frac{3p}{4}.
\]
All three orthogonal directions have eigenvalue
\[
\lambda_2=\lambda_3=\lambda_4=\frac p4.
\]

Hence
\[
\gamma_{AB}(p)
=
\left(1-\frac{3p}{4}\right)^2 + 3\left(\frac p4\right)^2.
\]
Expanding,
\[
\gamma_{AB}(p)
=
1-\frac{3}{2}p+\frac{3}{4}p^2.
\]

This is the analytical formula used in the experiment:
\[
\gamma_{AB}(p)=1-\frac{3}{2}p+\frac{3}{4}p^2.
\]

\subsection*{Reduced density operators and reduced purities}

We now compute the reduced states of subsystems \(A\) and \(B\).

For the Bell state,
\[
\operatorname{Tr}_B(\rho_0)=\operatorname{Tr}_A(\rho_0)=\frac{I_2}{2}.
\]
Moreover,
\[
\operatorname{Tr}_B\!\left(\frac{I_4}{4}\right)=\frac{I_2}{2},
\qquad
\operatorname{Tr}_A\!\left(\frac{I_4}{4}\right)=\frac{I_2}{2}.
\]

Therefore
\[
\rho_A(p)=\operatorname{Tr}_B(\rho(p))
=
(1-p)\frac{I_2}{2}+p\frac{I_2}{2}
=
\frac{I_2}{2},
\]
and similarly
\[
\rho_B(p)=\frac{I_2}{2}.
\]

Thus the reduced states are independent of \(p\), and their purities are
\[
\gamma_A(p)=\operatorname{Tr}(\rho_A(p)^2)
=
\operatorname{Tr}\left[\left(\frac{I_2}{2}\right)^2\right]
=
\operatorname{Tr}\left(\frac{I_2}{4}\right)
=
\frac{1}{2},
\]
\[
\gamma_B(p)=\frac{1}{2}.
\]

Hence
\[
\gamma_A(p)=\gamma_B(p)=\frac12.
\]

\subsection*{Fidelity with respect to \texorpdfstring{$|\Psi^+\rangle$}{|Psi+>}}

Since the reference state is pure, the fidelity reduces to
\[
F_{\Psi^+}(p)
=
\langle\Psi^+|\rho(p)|\Psi^+\rangle.
\]

Substituting the expression for \(\rho(p)\),
\[
F_{\Psi^+}(p)
=
(1-p)\langle\Psi^+|\rho_0|\Psi^+\rangle
+
\frac p4 \langle\Psi^+|I_4|\Psi^+\rangle.
\]
Now
\[
\langle\Psi^+|\rho_0|\Psi^+\rangle = 1,
\qquad
\langle\Psi^+|I_4|\Psi^+\rangle = 1,
\]
so
\[
F_{\Psi^+}(p)
=
(1-p)+\frac p4
=
1-\frac{3p}{4}.
\]

Therefore
\[
F_{\Psi^+}(p)=1-\frac{3}{4}p.
\]

\subsection*{Concurrence}

The family \(\rho(p)\) is a Bell-diagonal Werner-type family, and its concurrence is
\[
C(p)=\max\left(0,\,2\lambda_{\max}-1\right),
\]
where \(\lambda_{\max}\) is the largest Bell-basis eigenvalue. In the present case,
\[
\lambda_{\max}=1-\frac{3p}{4},
\]
hence
\[
C(p)
=
\max\left(0,\,2\left(1-\frac{3p}{4}\right)-1\right)
=
\max\left(0,\,1-\frac{3p}{2}\right).
\]

Thus
\[
C(p)=\max\left(0,\,1-\frac{3}{2}p\right).
\]

The concurrence vanishes when
\[
1-\frac{3p}{2}=0
\quad\Longrightarrow\quad
p=\frac{2}{3}.
\]
This identifies the separability threshold for the depolarized family.

\subsection*{Pauli correlations}

We now compute the aligned Pauli correlations
\[
c_{ii}(p)=\operatorname{Tr}\bigl[\rho(p)(\sigma_i\otimes\sigma_i)\bigr],
\qquad i\in\{x,y,z\}.
\]

Using the linearity of the trace,
\[
c_{ii}(p)
=
(1-p)\operatorname{Tr}\bigl[\rho_0(\sigma_i\otimes\sigma_i)\bigr]
+
\frac p4 \operatorname{Tr}\bigl[I_4(\sigma_i\otimes\sigma_i)\bigr].
\]

Since \(\sigma_i\otimes\sigma_i\) is traceless,
\[
\operatorname{Tr}\bigl[I_4(\sigma_i\otimes\sigma_i)\bigr]
=
\operatorname{Tr}(\sigma_i\otimes\sigma_i)=0.
\]
Therefore
\[
c_{ii}(p)=(1-p)c_{ii}(0).
\]

For the Bell state \(|\Psi^+\rangle\), one has
\[
\langle\Psi^+|\sigma_x\otimes\sigma_x|\Psi^+\rangle = 1,
\]
\[
\langle\Psi^+|\sigma_y\otimes\sigma_y|\Psi^+\rangle = 1,
\]
\[
\langle\Psi^+|\sigma_z\otimes\sigma_z|\Psi^+\rangle = -1.
\]

Hence
\[
c_{xx}(p)=1-p,
\qquad
c_{yy}(p)=1-p,
\qquad
c_{zz}(p)=-(1-p).
\]


\subsection*{Summary of analytical expressions}

Collecting all results, the depolarized Bell-state family satisfies
\[
\boxed{
\begin{aligned}
\gamma_{AB}(p) &= 1-\frac{3}{2}p+\frac{3}{4}p^2, \\
\gamma_A(p) &= \gamma_B(p)=\frac12, \\
F_{\Psi^+}(p) &= 1-\frac34 p, \\
C(p) &= \max\left(0,\,1-\frac32 p\right), \\
c_{xx}(p) &= 1-p,\qquad c_{yy}(p)=1-p,\qquad c_{zz}(p)=-(1-p).
\end{aligned}
}
\]

These are exactly the expressions implemented in the analytical branch of Experiment 3 and used as reference for numerical validation.

%========================================================
%========================================================
%========================================================
\section{References}

\begin{itemize}
    \item M. A. Nielsen and I. L. Chuang, \textit{Quantum Computation and Quantum Information}, Cambridge University Press.
\end{itemize}

\end{document}